\section{Example workflows}
\label{sec:examples}

Every listing below is an excerpt of a function in
\code{paper/phynetpy-1.0.0/listings/paper_listings.py}, which runs as a script
and fails loudly if any listing stops matching the API. Figure~\ref{fig:workflow}
shows how the pieces fit together.

\begin{figure}[t]
  \centering
  \includegraphics[width=0.78\linewidth]{workflow.pdf}
  \caption{A simulate--infer--evaluate loop in the public API. Data enters
  either by reading a file or by simulating from a known network; the two verbs
  act on it; the evaluation step closes the loop. Because \code{simulate}
  returns the same kind of object that \code{infer} consumes, a recovery check
  needs no intermediate conversion.}
  \label{fig:workflow}
\end{figure}

\subsection{Building and inspecting a network}

Reticulations are written in extended Newick and the resulting object supports
the structural queries directly.

\begin{lstlisting}
from phynetpy import blobs, count_reticulations, level, read_newick

net = read_newick("((A:0.1,(B:0.2)#H1:0.1):0.3,(#H1:0.2,C:0.3):0.2);")

count_reticulations(net)   # 1
level(net)                 # 1
net.is_acyclic()           # True
len(blobs(net))            # 4
\end{lstlisting}

The same type can be assembled node by node, which is what a method that
proposes topologies does internally:

\begin{lstlisting}
from phynetpy import Edge, Network, Node

root, internal = Node("Root"), Node("I1")
a, b, c = Node("A"), Node("B"), Node("C")
net = Network()
net.add_nodes([root, internal, a, b, c])
for edge in (Edge(root, internal), Edge(root, c),
             Edge(internal, a), Edge(internal, b)):
    net.add_edges(edge)
\end{lstlisting}

\subsection{Scoring a network, then inferring one}

This listing shows the unit discipline in use. The simulator writes branch
lengths in expected substitutions per site; the gene-tree criteria read
coalescent units. Converting is a call, not an assumption.

\begin{lstlisting}
from phynetpy import BranchLengthUnit, convert_network_branch_lengths
from phynetpy.criteria import Likelihood, PseudoLikelihood
from phynetpy.infer import infer, score, simulate
from phynetpy.models import MSC

theta = 0.5
gts = simulate(MSC(theta=theta), taxa=6, n=200, data="gene_trees", seed=3)
truth = convert_network_branch_lengths(
    gts.true_network, BranchLengthUnit.COALESCENT_2N, theta=theta
)

score(truth, gts, model=MSC(theta=theta), criterion=Likelihood())
#  -514.956
score(truth, gts, model=MSC(theta=theta), criterion=PseudoLikelihood())
# -1883.194

result = infer(gts, criterion=PseudoLikelihood(), max_reticulations=0)
result.method        # 'InferNetwork_MPL'
result.score         # -1882.494
\end{lstlisting}

Changing \code{criterion} to \code{Bayesian()} turns the same call into a
posterior sample, and \code{result.posterior} is then populated. The rest of the
call, and the type it returns, do not change.

\subsection{Simulation and recovery}

Because \code{simulate} produces a data-axis object, a recovery check is two
lines. This is the pattern behind the benchmark in
Section~\ref{sec:headtohead} and behind \code{Examples/sim_recovery.py}.

\begin{lstlisting}
from phynetpy.GraphUtils import mu_distance
from phynetpy.infer import infer, simulate

sim = simulate(MSC(theta=0.5), taxa=6, n=200, data="gene_trees", seed=11)
recovered = infer(sim, criterion=PseudoLikelihood(), max_reticulations=0)

mu_distance(sim.true_network, recovered.best)   # 0
\end{lstlisting}

The generating network is attached to the simulated data as
\code{.true_network}, so the comparison does not require tracking it separately.

\subsection{Inspecting and extending the method set}

The set of available methods is queryable rather than documented separately:

\begin{lstlisting}
from phynetpy.infer import validity_matrix

for data, row in validity_matrix().items():
    print(data, row)
\end{lstlisting}

\begin{lstlisting}[language={}, basicstyle=\ttfamily\scriptsize]
GeneTrees        {'MDC': '-', 'Likelihood': 'InferNetwork_ML',
                  'PseudoLikelihood': 'InferNetwork_MPL', 'Bayesian': 'MCMC_GT'}
Alignment        {'MDC': 'x', 'Likelihood': '-',
                  'PseudoLikelihood': 'x', 'Bayesian': 'MCMC_SEQ'}
BiallelicMarkers {'MDC': 'x', 'Likelihood': 'MLE_BiMarkers',
                  'PseudoLikelihood': '-', 'Bayesian': 'MCMC_BiMarkers'}
\end{lstlisting}

Adding a criterion means declaring what it accepts and registering an engine for
it. The objective below is deliberately trivial -- total branch length -- so that
the mechanism rather than the statistics is visible:

\begin{lstlisting}
from phynetpy.criteria import Criterion
from phynetpy.data import GeneTrees
from phynetpy.infer import Engine, register, score
from phynetpy.models import MSC

class TreeLength(Criterion):
    accepts_data = (GeneTrees,)
    scorable = True
    use_branch_lengths = False

@register(GeneTrees, MSC, TreeLength)
class TreeLengthEngine(Engine):
    method = "TreeLength"

    def infer(self, data):
        raise NotImplementedError("scoring only")

    def score(self, network, data, optimize=False):
        return sum(e.get_length() or 0.0 for e in network.E())
\end{lstlisting}

After that, \code{score(net, data, criterion=TreeLength())} dispatches to the new
engine, \code{validity_matrix()} shows the new column, and an unregistered
combination involving \code{TreeLength} raises \code{NotImplementedError} rather
than silently falling back. A real criterion would additionally implement
\code{infer}, and would get the search drivers, topology moves and parameter
optimiser from the shared layer rather than reimplementing them.
